\chapter{The Living Reactor: Milk Fermentation and the Screening of Strains}
\label{ch:ferment}

\begin{quote}\itshape
``Fermentation is the consequence of life without air.'' --- Louis Pasteur.
\end{quote}

\noindent
The unit operations of the previous chapters are, for all their mathematical
richness, \emph{clean}. A chromatographic column or an ultrafiltration cassette
does exactly what its conservation laws say it will, and the quantities we care
about---concentration, pressure, flux---can be measured more or less directly.
This chapter turns to an operation of a wholly different character: a
\emph{fermentation}, in which living bacteria, not a fixed packing, do the work.
The reactor is a beaker of milk; the catalyst is a community of lactic-acid
bacteria; the product is yogurt, or one of its many cousins. Three features make
this a separate intellectual problem, and organise the chapter.

First, the system is \emph{indirectly observed}. We cannot watch the bacteria
grow or count the molecules of lactose they consume; the one signal a small-scale
experiment yields cheaply and continuously is the \emph{pH}, a single curve that
folds together biomass, substrate, acid, and flavour. Section~\ref{sec:ferm-lab}
describes how the real experiment is run and read.

Second, the system is \emph{irreducibly stochastic}. ``Identical'' batches
differ---because the inoculum, the milk lot, and the incubator are never quite the
same---and a single batch wanders, because growth is a noisy process.
Sections~\ref{sec:ferm-model} and \ref{sec:ferm-noise} build a mechanistic model
that nonetheless wears its randomness on its sleeve, separating \emph{variability}
(batch-to-batch biology) from \emph{uncertainty} (the measurement).

Third, the design variable is \emph{combinatorial}. The central question of the
dairy fermentation industry is not ``what temperature?'' but ``\emph{which
strains?}''---typically a blend of two to five, drawn from a library of dozens of
candidates, chosen so that a new culture reproduces the product of an old one
being retired. Sections~\ref{sec:ferm-screen} and \ref{sec:ferm-analysis} develop
the design and the analysis for exactly this: \emph{covering arrays} to choose the
combinations, and \emph{tree-ensemble} learning to find, from noisy data, which
strains matter.

%==============================================================================
\section{The real-world experiment}
\label{sec:ferm-lab}

A yogurt fermentation is deceptively simple to perform and surprisingly hard to
read. Milk is pasteurised, cooled to an incubation temperature (typically
$40$--$45^\circ$C for a thermophilic yogurt culture), and \emph{inoculated} with a
small dose of a \emph{starter}---a defined mixture of bacterial strains. The
bacteria ferment the milk sugar, \emph{lactose}, into \emph{lactic acid}; the acid
lowers the pH from that of fresh milk ($\approx 6.6$) and, as the casein micelles
lose their charge, the milk \emph{gels} into the familiar set. The fermentation is
stopped---by cooling---when a target pH is reached, conventionally $\mathrm{pH}\,
4.6$, the isoelectric neighbourhood of casein where the gel is firm.

\subsection{What is measured, and what is not}

The workhorse measurement is a \emph{pH time series}, $\{\mathrm{pH}(t_i)\}$,
logged by a probe every few minutes over the several hours of the run. From it the
technologist reads a handful of features (Figure~\ref{fig:ferm-acid}, left): the
\emph{time to gelation} (pH $5.2$), the \emph{time to set} (pH $4.6$)---the single
most-reported number in strain screening---the \emph{maximum acidification rate}
$\max_t(-\dv{\mathrm{pH}}{t})$, and the \emph{post-acidification}, the further
drop during cold storage that, in excess, is a defect. A complementary off-line
assay is the \emph{titratable acidity}, the amount of base needed to bring a
sample back to neutrality, reported in $^\circ$Dornic or as \% lactic acid; it
measures total acid where pH measures only free protons, the two differing because
milk is a \emph{buffer}.

What is \emph{not} measured, routinely, is everything mechanistic: the viable cell
count of each strain (a laborious plating assay), the residual lactose, the
concentration of each acid and of the aroma compounds (chromatography). The pH
curve is thus an \emph{indirect, aggregate} indicator---a shadow cast by several
latent variables at once. Any model of this system is, in effect, a theory of what
that shadow implies.

\begin{databox}{the pH curve and its features}
The calibration target for the model is not a single number but the \emph{shape}
of $\mathrm{pH}(t)$: a fresh-milk plateau, a lag, a steep acidification through
the gelation region, and a levelling-off as the culture exhausts itself. The
features above (set time, max rate, post-acidification) are the quantities against
which strains are compared and against which a model is judged realistic. In the
package they are computed by \code{fermentation.fingerprint} and
\code{time\_to\_ph}.
\end{databox}

\subsection{Why strains, and the replacement problem}

The flavour, texture, acidity, and shelf-life of a cultured product are set
overwhelmingly by \emph{which bacteria} ferment it. Classical yogurt is the
symbiosis of two: \emph{Streptococcus thermophilus} and \emph{Lactobacillus
delbrueckii} subsp.\ \emph{bulgaricus}. Adjunct \emph{probiotic} species
(\emph{L.\ acidophilus}, \emph{Bifidobacterium}) are added for health claims.
Strains within a species differ markedly in growth rate, acid tolerance,
acidifying power, and aroma.

The recurring industrial problem is \emph{replacement}: a strain in a marketed
product is discontinued (a supplier change, a phage sensitivity, a regulatory
issue), and a new strain or blend must be found that reproduces the old product's
sensory and kinetic profile. This is a \emph{matching} problem against a moving
target, over a combinatorial space of blends---and it is the motivation for the
designs and analyses of this chapter.

\begin{appbox}{cultured dairy and beyond}
The same model and methods describe set and stirred yogurt, drinking yogurt,
cultured buttermilk, sour cream, fresh cheeses, and the acidification phase of
hard-cheese make. More broadly, any \emph{batch microbial fermentation read out by
a single aggregate signal}---a probiotic ferment, a sourdough, a silage, even a
bioprocess monitored only by off-gas or pH---shares the same structure: latent
biology, an indirect observable, and a combinatorial choice of inoculum.
\end{appbox}

%==============================================================================
\section{A mechanistic--stochastic model}
\label{sec:ferm-model}

We seek the \emph{simplest model that retains the right physics}---the guiding
principle of this whole package. There is no single accepted mechanistic model of
yogurt fermentation as there is for chromatography; what follows assembles
standard, well-validated pieces from predictive microbiology and bioreaction
engineering into a generator whose output behaves, qualitatively and
quantitatively, like a real fermentation, and whose randomness is explicit.

\subsection{State variables}

For a consortium of $n$ strains the state at time $t$ (in hours) is
%
\begin{equation}
\bm y(t) = \bigl(\, X_1,\dots,X_n,\; Q_1,\dots,Q_n,\; S,\; L,\; A \,\bigr),
\end{equation}
%
where $X_i$ is the biomass of strain $i$ (a dimensionless scaled cell density),
$Q_i$ is its \emph{physiological state} (the lag variable introduced below), $S$
is the lactose concentration (g/L), $L$ is the lactic-acid concentration
(mmol/L), and $A$ is an aroma proxy (acetaldehyde, in arbitrary units). The pH is
deliberately \emph{not} a state variable: it is recovered from $L$ at every
instant through the milk titration curve of Section~\ref{sec:ferm-ph}, so that the
acid the bacteria produce feeds back, without delay, into their own inhibition.

\subsection{Growth: the Baranyi--Roberts law}

Each strain grows according to the model of Baranyi and Roberts~\cite{baranyi},
the standard of predictive microbiology, because it represents the \emph{lag
phase} mechanistically rather than by a switch. Its idea is that cells emerging
from a dormant inoculum must first build up some limiting intracellular pool before
they can divide; an \emph{adjustment function} $\alpha_i(t)\in[0,1]$ gates the
growth rate from zero up to one as that pool fills,
%
\begin{equation}
\alpha_i(t) = \frac{Q_i(t)}{1+Q_i(t)},
\qquad
\dv{Q_i}{t} = \mu_{\max,i}\,\gamma_{T,i}\,Q_i,
\label{eq:baranyi-alpha}
\end{equation}
%
so that $Q_i$ grows exponentially from its inoculation value $Q_i(0)=Q_{0,i}$ and
$\alpha_i$ follows a smooth logistic switch. Integrating
Eq.~\eqref{eq:baranyi-alpha} gives the lag duration in closed form,
%
\begin{equation}
\lambda_i = \frac{\ln\!\bigl(1 + 1/Q_{0,i}\bigr)}{\mu_{\max,i}\,\gamma_{T,i}},
\label{eq:lag}
\end{equation}
%
which exposes the modelling pay-off: a \emph{single} parameter $Q_{0,i}$ sets the
lag, and---because the physiological state of a real inoculum is genuinely
variable---making $Q_{0,i}$ random (Section~\ref{sec:ferm-noise}) makes the lag
random, the most conspicuous source of run-to-run spread in practice. Biomass then
follows a lag-gated, substrate-limited, acid-inhibited logistic law,
%
\begin{equation}
\dv{X_i}{t}
= \alpha_i(t)\,\underbrace{\mu_{\max,i}\,
   \gamma_{T,i}(T)\,
   \frac{S}{K_{S,i}+S}\,
   I_i(\mathrm{pH})\,
   \Gamma_i(\bm X)}_{\textstyle \mu_i \;=\; \text{specific growth rate}}
\;X_i\,
\Bigl(1-\frac{X_{\mathrm{tot}}}{X_{\max,i}}\Bigr),
\label{eq:growth}
\end{equation}
%
with $X_{\mathrm{tot}}=\sum_j X_j$. The four dimensionless factors building $\mu_i$
are the four sub-models of the next paragraphs: a temperature factor
$\gamma_{T,i}$, a Monod substrate factor, an acid-inhibition factor $I_i$, and an
interaction factor $\Gamma_i$. The Monod term $S/(K_{S,i}+S)$~\cite{monod} is the
classical saturating dependence on the limiting nutrient, with half-saturation
$K_{S,i}$; the logistic factor caps the total population at the carrying capacity.

\subsection{Temperature: the cardinal model}

Bacterial growth rate depends on temperature in a strongly asymmetric way: it
rises gently from a minimum, peaks at an optimum, then collapses over a few degrees
to a maximum above which growth ceases. The \emph{cardinal temperature model with
inflection} of Rosso et al.~\cite{rosso} captures this with three biologically
meaningful parameters---the cardinal temperatures $T_{\min}$, $T_{\mathrm{opt}}$,
$T_{\max}$---and no others:
%
\begin{equation}
\gamma_{T}(T) =
\frac{(T-T_{\max})(T-T_{\min})^2}
     {(T_{\mathrm{opt}}-T_{\min})
      \bigl[(T_{\mathrm{opt}}-T_{\min})(T-T_{\mathrm{opt}})
           -(T_{\mathrm{opt}}-T_{\max})(T_{\mathrm{opt}}+T_{\min}-2T)\bigr]},
\label{eq:rosso}
\end{equation}
%
clamped to $\gamma_T=0$ outside $[T_{\min},T_{\max}]$ and normalised so that
$\gamma_T(T_{\mathrm{opt}})=1$. Multiplying $\mu_{\max}$ by
Eq.~\eqref{eq:rosso} turns incubation temperature into a genuine process factor and
gives each strain a distinct thermal niche (Figure~\ref{fig:ferm-sub}, left):
\emph{S.\ thermophilus} peaks near $42^\circ$C, \emph{L.\ bulgaricus} near
$45^\circ$C, so the blend's behaviour shifts with temperature in a way no single
strain reproduces.

\subsection{Acid inhibition}

As acid accumulates the falling pH eventually arrests growth, and it does so at a
strain-specific threshold---the essence of \emph{acid tolerance}. We use a linear
ramp from full activity at a reference pH down to zero at the strain's minimum
growth pH $\mathrm{pH}_{\min,i}$,
%
\begin{equation}
I_i(\mathrm{pH}) =
\operatorname{clamp}\!\left(
\frac{\mathrm{pH}-\mathrm{pH}_{\min,i}}{\mathrm{pH}_{\mathrm{ref}}-\mathrm{pH}_{\min,i}},
\;0,\;1\right).
\label{eq:acid-inhib}
\end{equation}
%
The consequences are exactly those seen in the dairy: \emph{S.\ thermophilus}, with
a high $\mathrm{pH}_{\min}\approx 4.7$, stalls early and cannot set milk on its
own; \emph{L.\ bulgaricus}, tolerant to $\approx 3.8$, drives the final
acidification and, through maintenance metabolism, the post-acidification.

\subsection{Strain interactions: proto-cooperation}

The yogurt symbiosis is the textbook case of \emph{proto-cooperation}: \emph{S.\
thermophilus} releases formate and carbon dioxide that stimulate \emph{L.\
bulgaricus}, whose proteolysis in turn liberates the peptides and amino acids that
\emph{S.\ thermophilus} needs. We model this as a saturating, pairwise
multiplicative boost to the growth rate,
%
\begin{equation}
\Gamma_i(\bm X) = \prod_{j\neq i}
\left(1 + k_{ij}\,\frac{X_j}{X_j+K_c}\right),
\label{eq:interaction}
\end{equation}
%
where the entry $k_{ij}$ of the consortium's \emph{interaction matrix} is the
strength with which strain $j$ affects strain $i$: positive for cooperation,
negative (with $k_{ij}>-1$) for antagonism, zero for independence (the default).
The half-saturation $K_c$ sets the biomass at which a partner's influence is half
its maximum. The effect is decisive (Figure~\ref{fig:ferm-coop}): with
$k_{ij}=0$ the same two strains, blended, stall above the set point; switch
cooperation on and the pair sets the milk together and faster than either alone.

\subsection{Acid and aroma production}

Acid is produced by both growth and maintenance, the \emph{Luedeking--Piret}
relation~\cite{luedeking} of bioreaction engineering:
%
\begin{equation}
\dv{L}{t} =
\sum_i \Bigl( a_i\,\dv{X_i}{t} \;+\; b_i\,I_i(\mathrm{pH})\,X_i \Bigr),
\label{eq:luedeking}
\end{equation}
%
with a growth-associated coefficient $a_i$ (mmol acid per unit biomass formed) and
a non-growth maintenance coefficient $b_i$ (mmol acid per unit biomass per hour).
The maintenance term is gated by the same acid-inhibition factor
$I_i$ of Eq.~\eqref{eq:acid-inhib}: a strain that has reached its minimum pH stops
metabolising, which prevents an acid-sensitive strain from acidifying past its own
stall point yet lets a tolerant strain continue, reproducing post-acidification.
Lactose is consumed stoichiometrically through the growth-associated yield,
$\dv{S}{t} = -\sum_i Y_i^{-1}\,\dv{X_i}{t}$ (gated to vanish once $S$ is
exhausted), and the aroma proxy follows $\dv{A}{t} = \sum_i c_i\,\dv{X_i}{t} -
k_{\mathrm{dec}}A$, so that the flavour fingerprint shifts with the community
composition.

\subsection{From acid to pH: milk as a buffer}
\label{sec:ferm-ph}

The link from the acid the bacteria make to the pH we measure is milk's
\emph{buffering capacity}. Casein, phosphate, and citrate resist the pH change, so
the titration curve is a shallow sigmoid, not the straight line a strong acid in
water would give. Rather than model each buffering species we use an empirical
titration curve that reproduces the measured shape,
%
\begin{equation}
\mathrm{pH}(L) = \mathrm{pH}_\infty +
\frac{\mathrm{pH}_0 - \mathrm{pH}_\infty}{1 + (L/L_{50})^{n_{\mathrm{buf}}}},
\label{eq:titration}
\end{equation}
%
where $\mathrm{pH}_0\approx 6.6$ is fresh milk, $\mathrm{pH}_\infty$ the floor the
pH asymptotes toward, $L_{50}$ the acid at the curve's midpoint---the practical
handle on buffering capacity---and $n_{\mathrm{buf}}$ the steepness through the
gelation region. Because $L_{50}$ and the initial lactose $S_0$ vary appreciably
between milk lots, they are the natural carriers of one part of the batch
variability (Section~\ref{sec:ferm-noise}), and they live on the model's
\code{Milk} object. Figure~\ref{fig:ferm-sub} (right) shows how a lower-buffer
milk lot reaches any given pH at less acid.

\subsection{The complete system}

Collecting the pieces, the deterministic core is the $2n+3$ coupled ordinary
differential equations \eqref{eq:baranyi-alpha}--\eqref{eq:titration}, integrated
(when noise is off) by a stiff solver (\code{scipy}'s LSODA, which handles the
abrupt lag-to-growth transition). Figure~\ref{fig:ferm-acid} shows a representative
solution: the single measurable pH curve on the left, and on the right the four
latent trajectories---two biomasses, lactose, acid---that the pH curve silently
encodes. This is the heart of the matter: the experiment sees only the left panel
and must infer the right.

\begin{figure}[t]
\centering
\includegraphics[width=\linewidth]{ferm_acidification.png}
\caption[A simulated yogurt fermentation]{A representative yogurt batch (a $50{:}50$
\emph{S.\ thermophilus}/\emph{L.\ bulgaricus} blend at $43^\circ$C). \emph{Left:}
the only observable---the pH curve---with the gelation and set-point milestones.
\emph{Right:} the latent states the pH curve aggregates and which the experiment
never measures directly: the two strains' biomass, the consumed lactose, and the
accumulated lactic acid.}
\label{fig:ferm-acid}
\end{figure}

\begin{figure}[t]
\centering
\includegraphics[width=\linewidth]{ferm_submodels.png}
\caption[The three sub-models]{The three constitutive sub-models. \emph{Left:} the
Rosso cardinal-temperature factor $\gamma_T(T)$, Eq.~\eqref{eq:rosso}, giving each
strain an asymmetric thermal niche. \emph{Centre:} the Baranyi adjustment function
$\alpha(t)$, Eq.~\eqref{eq:baranyi-alpha}, whose single parameter $Q_0$ sets the
lag. \emph{Right:} the milk titration curve $\mathrm{pH}(L)$,
Eq.~\eqref{eq:titration}, for three buffering capacities.}
\label{fig:ferm-sub}
\end{figure}

\begin{figure}[t]
\centering
\includegraphics[width=0.8\linewidth]{ferm_cooperation.png}
\caption[Strain personalities and the symbiosis]{Strain personalities. \emph{S.\
thermophilus} alone acidifies fast but stalls above the set point (high
$\mathrm{pH}_{\min}$); \emph{L.\ bulgaricus} alone is slow but acid-tolerant. The
cooperative blend (solid green) sets fastest, while the \emph{same} blend with the
interaction switched off (dashed grey) barely sets---the proto-cooperation of
Eq.~\eqref{eq:interaction} is what makes yogurt.}
\label{fig:ferm-coop}
\end{figure}

\begin{limitbox}{the fermentation model}
\begin{itemize}
  \item \textbf{Phenomenological, not first-principles.} Unlike the chromatography
  models, this is an assembly of standard empirical laws chosen to reproduce the
  \emph{shape} and \emph{sensitivities} of real fermentations, not a derivation
  from molecular biology. Its parameters are illustrative, not measured from a
  specific organism.
  \item \textbf{The titration curve is empirical.} Eq.~\eqref{eq:titration}
  mimics milk buffering without modelling the individual buffering species; it is
  valid over the normal fermentation range, not for extreme acidification.
  \item \textbf{Well-mixed and isothermal.} Spatial gradients (a real set gel is
  not stirred) and the small self-heating of metabolism are ignored.
  \item \textbf{Aroma is a single proxy.} Real flavour is dozens of compounds;
  the model carries one, as a stand-in for the sensory dimension.
\end{itemize}
\end{limitbox}

%==============================================================================
\section{Three faces of randomness}
\label{sec:ferm-noise}

The request that gave rise to this model was to ``represent \emph{all} the
uncertainty.'' It is essential to keep distinct three kinds, which enter at
different places and must be estimated differently
(Figure~\ref{fig:ferm-noise}).

\paragraph{Batch variability (aleatoric).}
Nominally identical batches differ because their \emph{parameters} differ. We model
this hierarchically: a population mean $\bar\theta$ for each batch parameter, and a
per-batch draw with a lognormal multiplier,
%
\begin{equation}
\theta^{(b)} = \bar\theta \cdot \exp\!\bigl(\sigma_\theta\,Z\bigr),
\qquad Z\sim\mathcal N(0,1),
\label{eq:lognormal-jitter}
\end{equation}
%
so that $\sigma_\theta$ is, for small values, the fractional coefficient of
variation, and the multiplier is always positive. The jittered parameters are the
inoculum size, the lag states $Q_{0,i}$, the maximum growth rates, the milk
buffering $L_{50}$ and lactose $S_0$, and an additive incubator-temperature offset.
This is the irreducible spread \emph{between} batches; it does not shrink with a
better instrument, only with better process control. It is generated by
\code{sample\_batch} (built on the project's
\code{perturbation.jitter\_parameters}).

\paragraph{Process noise (within-batch).}
A single fermentation also \emph{wanders}: growth is a stochastic birth process,
and local fluctuations accumulate. We add an It\^o diffusion term to the biomass
equations, turning Eq.~\eqref{eq:growth} into the stochastic differential equation
%
\begin{equation}
\dd X_i = f_i(\bm y)\,\dd t + \sigma_{\mathrm{proc}}\,X_i\,\dd W_i,
\label{eq:sde}
\end{equation}
%
with $f_i$ the deterministic drift and $W_i$ independent Wiener processes; the
multiplicative form keeps the noise proportional to the population and biomass
non-negative. It is integrated by the Euler--Maruyama scheme on a fine grid,
%
\begin{equation}
X_i^{(k+1)} = X_i^{(k)} + f_i\,\Delta t
 + \sigma_{\mathrm{proc}}\,X_i^{(k)}\,\sqrt{\Delta t}\;Z_i^{(k)},
 \qquad Z_i^{(k)}\sim\mathcal N(0,1),
\label{eq:euler-maruyama}
\end{equation}
%
with each updated state floored at zero. Setting $\sigma_{\mathrm{proc}}=0$
recovers the deterministic solver exactly; switching it on makes individual
trajectories \emph{wobble} rather than being clean shifts of one curve.

\paragraph{Measurement uncertainty (epistemic).}
Finally, the pH probe is imperfect. We sample the latent pH curve at the probe's
(coarser) times and add the project's standard measurement model---additive
Gaussian noise plus an optional calibration bias,
%
\begin{equation}
\widetilde{\mathrm{pH}}(t_i)
 = \mathrm{pH}(t_i) + \eps_i + \beta,
\qquad \eps_i\sim\mathcal N(0,\sigma_{\mathrm{meas}}^2),
\end{equation}
%
through \code{observe\_ph}, which reuses \code{perturbation.add\_measurement\_noise}
so the fermentation probe is treated identically to every other virtual
instrument. This noise \emph{does} shrink with a better instrument, and it can in
principle be averaged away by replicate measurements of one batch---which is
exactly what distinguishes it from the variability above.

\begin{figure}[t]
\centering
\includegraphics[width=\linewidth]{ferm_uncertainty.png}
\caption[Three kinds of randomness]{The three faces of randomness, kept separate.
\emph{Left:} measurement uncertainty---probe noise on a single batch. \emph{Centre:}
batch variability---the spread of forty whole batches drawn from the population of
Eq.~\eqref{eq:lognormal-jitter}. \emph{Right:} process noise---eight runs of the
SDE, Eq.~\eqref{eq:sde}, each wandering on its own path. Telling these apart is the
whole point of separating uncertainty from variability.}
\label{fig:ferm-noise}
\end{figure}

\subsection{The product fingerprint and the replacement objective}
\label{sec:ferm-fingerprint}

To make the replacement problem of Section~\ref{sec:ferm-lab} quantitative we
summarise a batch by a \emph{fingerprint} vector $\bm f$---the gelation and set
times, the final pH, the post-acidification, the maximum rate and its timing, the
final aroma, and the community composition---computed by \code{fingerprint}. Two
products are then compared by a weighted, scaled distance to a reference $\bm r$
(the incumbent strain being retired),
%
\begin{equation}
D(\bm f,\bm r) = \sqrt{\;\sum_k w_k
\left(\frac{f_k - r_k}{s_k}\right)^{\!2}\,},
\label{eq:fingerprint-dist}
\end{equation}
%
with characteristic scales $s_k$ to non-dimensionalise features of different units
and weights $w_k$ to emphasise the attributes that matter for a given product. A
milestone that a candidate never reaches (a set point never attained) is penalised
by one full scale-unit. Minimising $D$ over the blend composition and process
conditions \emph{is} the replacement problem, and---being a scalar objective over a
design space---it plugs directly into the Bayesian-optimisation machinery of
Chapter~\ref{ch:design}.

%==============================================================================
\section{Designing the screen: covering arrays and their relatives}
\label{sec:ferm-screen}

We now confront the combinatorial design problem. Suppose a library of $k=50$
candidate strains, a budget of $N=200$ experiments, and the practical constraint
that each experiment is a blend of between two and five strains. Which blends
should we run?

\subsection{The combinatorial explosion and the pairwise idea}

Exhaustive testing is hopeless: the number of $2$-to-$5$-strain blends of $50$
strains is
%
\begin{equation}
\sum_{m=2}^{5}\binom{50}{m}
= \binom{50}{2}+\binom{50}{3}+\binom{50}{4}+\binom{50}{5}
\approx 2.37\times 10^{6}.
\end{equation}
%
But we do not need every blend. If the response is governed mainly by individual
strains and by \emph{pairwise} strain interactions---a reasonable first
hypothesis, and the same order of approximation a factorial design makes for
continuous factors---then it suffices that \emph{every pair of strains appears
together in at least one run}. There are only $\binom{50}{2}=1225$ such pairs, a
number commensurate with a $200$-run budget. This is the principle of a
\emph{strength-2 covering array}.

\subsection{Covering arrays}

A covering array $\mathrm{CA}(N;t,k,v)$ is an $N\times k$ array over $v$ symbols
such that in \emph{every} choice of $t$ columns, all $v^t$ combinations of symbols
appear in at least one row~\cite{colbourn}. The parameter $t$ is the
\emph{strength}. The classical use is software testing, where $k$ configuration
options each take $v$ values and a strength-$t$ array exercises every $t$-way
combination in far fewer than the $v^k$ exhaustive tests.

Our problem is the binary case $v=2$ (a strain is present or absent) at strength
$t=2$, but with a twist that the textbook construction does not address: a
\emph{block-size constraint}. A real fermentation blends a few strains, not
twenty-five, so each row must have between $g_{\min}=2$ and $g_{\max}=5$ ones. With
that constraint the object is equivalently a \emph{covering design}---a collection
of small blocks from a $k$-set such that every pair of points lies in some block.
A counting bound makes the budget concrete: a block of $g$ strains covers
$\binom{g}{2}$ pairs, so covering all $\binom{k}{2}$ pairs needs at least
%
\begin{equation}
N \;\ge\; \frac{\binom{k}{2}}{\binom{g_{\max}}{2}}
\label{eq:schonheim}
\end{equation}
%
runs (a Sch\"onheim-type lower bound)---here $1225/10 \approx 123$ runs of five,
so $200$ is comfortable, with redundancy to spare for the noise.

\paragraph{Greedy construction.}
The package builds the design greedily (\code{doe/covering.py}). It keeps the set
of still-uncovered pairs and, for each run, assembles a block by repeatedly adding
the strain that covers the most \emph{newly} uncovered pairs, seeding from a
rarely-used strain and breaking ties toward the least-used strain so that
appearances stay balanced (good replication for the later regression). The block
size is drawn across the full $[g_{\min},g_{\max}]$ range but \emph{biased toward
larger blocks while many pairs remain uncovered}---a size-five block covers ten
pairs, a size-two block only one---so coverage completes without sacrificing the
required size variety. Crucially the achieved coverage is \emph{reported}, not
assumed: \code{CoveringArrayDesign.coverage} returns the covered fraction, the
redundancy, and the appearance balance, so an under-powered budget is visible
rather than silently failing. Figure~\ref{fig:ferm-cover} shows coverage climbing
with the number of runs---faster for larger blocks---and the balanced strain usage
of the design actually used, which co-tests about $92\%$ of all $1225$ pairs with
each strain appearing roughly fifteen times.

\begin{figure}[t]
\centering
\includegraphics[width=\linewidth]{ferm_covering.png}
\caption[Covering-array coverage and balance]{\emph{Left:} the fraction of the
$1225$ strain pairs co-tested as runs accumulate, for three block-size caps; larger
blocks cover the space faster (Eq.~\eqref{eq:schonheim}). \emph{Right:} the design
actually used ($50$ strains, $200$ runs, blocks of $2$--$5$) spreads usage evenly
across strains---each appears $\approx 15$ times---which gives the downstream
analysis balanced replication.}
\label{fig:ferm-cover}
\end{figure}

\begin{prosbox}{covering arrays for screening}
\begin{itemize}
  \item \textbf{Combinatorial reach at linear cost.} Guarantees (near-)complete
  pairwise coverage of a huge combination space in a budget that scales with the
  number of pairs, not the number of combinations.
  \item \textbf{Honest diagnostics.} Coverage fraction, redundancy, and balance are
  computed, so the design's adequacy is measurable before a single experiment.
  \item \textbf{Balanced and model-agnostic.} Even strain usage suits any
  downstream analysis; nothing about the design presupposes a particular model.
\end{itemize}
\end{prosbox}

\begin{consbox}{covering arrays for screening}
\begin{itemize}
  \item \textbf{Pairwise by default.} Strength-2 targets main effects and two-way
  interactions; genuine three-way synergies need strength-3, at far higher cost.
  \item \textbf{Coverage, not estimability.} It guarantees that every pair is
  \emph{seen}, not that every effect is \emph{orthogonally estimable}; the analysis
  must cope with correlated, non-orthogonal columns.
  \item \textbf{Greedy, not optimal.} The construction is a good heuristic, not a
  provably minimum design; the block-size cap can leave a few pairs uncovered.
\end{itemize}
\end{consbox}

\subsection{Other designs for strain screening}

Covering arrays are one of several families that attack high-dimensional discrete
screening; a mature campaign chooses among them by what it assumes
(Table~\ref{tab:ferm-screen}).

\emph{Fractional factorial and Plackett--Burman designs}~\cite{schmidtbox,plackett}
treat each strain as a two-level ($\pm$) factor and use a carefully chosen
\emph{fraction} of the $2^k$ runs. A Plackett--Burman design, built from a Hadamard
matrix, estimates all $k$ main effects in as few as $k+1$ runs---supremely
efficient---but \emph{confounds} (aliases) main effects with two-way interactions,
so it is blind to the strain synergies that matter most in fermentation. Its
resolution-III economy is the mirror image of the covering array's interaction
awareness.

\emph{Supersaturated designs} go further, with \emph{fewer runs than factors}
($N<k$); they can only flag a sparse handful of dominant strains and rely on
effect sparsity. \emph{Group-testing} designs share the same spirit as the covering
array---pool many items per test and deconvolve---and are optimal when the goal is
to \emph{find a few active strains} rather than to map interactions.

When the response depends on the \emph{proportions} of strains rather than merely
their presence, the proper framework is \emph{mixture design}~\cite{cornell}: the
blend fractions lie on a simplex $\sum_i x_i = 1$, and \emph{simplex-lattice} or
\emph{simplex-centroid} designs place points to fit a Scheff\'e mixture polynomial
%
\begin{equation}
y = \sum_i \beta_i x_i + \sum_{i<j}\beta_{ij}x_i x_j + \cdots,
\label{eq:scheffe}
\end{equation}
%
which has no intercept because the fractions are not independent. Mixture and
combinatorial views are complementary: a covering array first screens \emph{which}
strains belong in the blend; a mixture design then optimises \emph{how much} of
each. Finally, \emph{$D$-optimal} designs choose runs to maximise $\det(\bm
X\transp\bm X)$ for a \emph{specified} model---ideal once the model form is known,
but circular for screening, where the model is what we are trying to discover.

\begin{table}[h]
\centering
\renewcommand{\arraystretch}{1.35}
\caption{Screening designs for a large strain library.}
\label{tab:ferm-screen}
\begin{tabular}{@{}p{0.22\linewidth} p{0.30\linewidth} p{0.20\linewidth}
                  p{0.18\linewidth}@{}}
\toprule
\textbf{Design} & \textbf{Targets} & \textbf{Runs} & \textbf{Blind to} \\
\midrule
Covering array (strength 2) & Presence + all pairwise co-occurrence &
$\sim\binom{k}{2}/\binom{g_{\max}}{2}$ & Three-way synergy \\
Plackett--Burman & Main effects only & $k+1$ & All interactions \\
Supersaturated & A few dominant strains & $<k$ & Most effects \\
Mixture (simplex) & Blend \emph{proportions} & Model-dependent &
Presence/absence framing \\
$D$-optimal & A pre-specified model & Model-dependent & Unmodelled terms \\
\bottomrule
\end{tabular}
\end{table}

%==============================================================================
\section{Analysing the data: from ANOVA to trees}
\label{sec:ferm-analysis}

The screen yields a $200\times 50$ binary \emph{presence matrix} $\bm X$ and a
noisy response $\bm y$---the final pH, or the (censored) set time. How do we learn
from it which strains, and which combinations, drive the outcome?

\subsection{Why the linear model struggles here}

The response-surface ANOVA of Chapter~\ref{ch:design} is the natural first thought,
but it fits this problem poorly. The design is \emph{not orthogonal} (the covering
array guarantees coverage, not the diagonal $\bm X\transp\bm X$ that makes
factorial effects independent); the structure is \emph{interaction-heavy and
threshold-like} (a strain's effect depends sharply on its acid-tolerant partners,
Eq.~\eqref{eq:interaction}); and to capture pairwise terms explicitly the linear
model would need $\binom{50}{2}=1225$ interaction columns against $200$ rows---a
hopelessly underdetermined fit. We need methods that discover interactions
\emph{implicitly} and tolerate correlated, high-dimensional inputs. Tree ensembles
do exactly this.

\subsection{Random forests}

A regression \emph{tree} (CART) recursively partitions the input space, at each
node choosing the split that most reduces the within-node variance of $y$; the
prediction in a leaf is the mean response there. A single deep tree captures
interactions automatically---a split on strain $j$ \emph{within} a branch already
conditioned on strain $i$ is precisely a two-way effect---but it is high-variance,
fitting noise. The \emph{random forest}~\cite{breiman} tames this by
\emph{bagging}: it grows $B$ trees, each on a bootstrap resample of the runs and
each node restricted to a random subset of the strains, and averages them,
%
\begin{equation}
\hat y(\bm x) = \frac1B \sum_{b=1}^{B} T_b(\bm x).
\end{equation}
%
Averaging decorrelated trees cuts variance without raising bias, yielding a robust
nonlinear predictor that needs little tuning---an excellent default for a noisy
screen.

\subsection{Gradient boosting and XGBoost}

\emph{Boosting} builds the ensemble the opposite way: instead of averaging
independent trees it adds them \emph{sequentially}, each new tree fitting the
errors the current ensemble still makes,
%
\begin{equation}
F_m(\bm x) = F_{m-1}(\bm x) + \nu\, h_m(\bm x),
\end{equation}
%
with $h_m$ a small tree and $\nu$ a learning-rate shrinkage. \code{XGBoost}~\cite{xgboost}
is the modern realisation, distinguished by a \emph{regularised} objective that it
optimises with a second-order Taylor expansion. Writing $g_i,h_i$ for the first and
second derivatives of the loss at run $i$, the objective for adding tree $f$ is
%
\begin{equation}
\mathcal L \simeq \sum_i \Bigl[g_i f(\bm x_i)
 + \tfrac12 h_i f(\bm x_i)^2\Bigr]
 + \underbrace{\gamma T + \tfrac12\lambda\|\bm w\|^2}_{\Omega(f)},
\label{eq:xgb}
\end{equation}
%
where the penalty $\Omega$ charges for the number of leaves $T$ and the magnitude
of the leaf weights $\bm w$, with $\gamma,\lambda$ controlling complexity. This
explicit regularisation, plus shrinkage and column subsampling, lets boosting
extract more signal than a forest when there is signal to extract, at the price of
more tuning and a greater risk of overfitting noise. Running both
(Figure~\ref{fig:ferm-imp}) is the honest practice: agreement between two very
different ensembles is strong evidence that a finding is real.

\subsection{Honest accuracy and importance}

Two diagnostics, computed for both models (\code{doe/importance.py}), make the
analysis trustworthy. The first is the \emph{cross-validated coefficient of
determination}: the data are split into $K$ folds, each predicted from a model
trained on the rest, and
%
\begin{equation}
R^2_{\mathrm{CV}} = 1 -
\frac{\sum_i \bigl(y_i - \hat y_{-k(i)}(\bm x_i)\bigr)^2}
     {\sum_i (y_i-\bar y)^2},
\end{equation}
%
where $\hat y_{-k(i)}$ is trained without the fold containing run $i$. Unlike the
in-sample $R^2$, which a flexible model can drive to one by memorising, $R^2_{\mathrm{CV}}$
measures genuine out-of-sample prediction.

The second is \emph{permutation importance}, a model-agnostic measure of how much
each strain matters: randomly shuffle the column for strain $j$ (destroying its
relationship with $y$ while preserving its marginal distribution), and record the
resulting \emph{drop} in score,
%
\begin{equation}
I_j = R^2(\text{model}) - \E_{\pi}\!\bigl[R^2(\text{model};\,
      \text{column } j \text{ permuted})\bigr].
\label{eq:permimp}
\end{equation}
%
A strain whose shuffling badly hurts prediction is important. Permutation
importance is used for \emph{both} ensembles---rather than each library's built-in
impurity or gain measure---precisely so the two are comparable and neither is
biased toward high-cardinality features.

\subsection{Does it work? Recovery from noise}

Because this is a virtual laboratory we can grade the analysis against ground
truth. Each library strain has an intrinsic acidifying power, which we can compute
by simulating it \emph{alone}---something the screen never does, since every run is
a $2$-to-$5$-strain blend. Figure~\ref{fig:ferm-imp} shows the payoff. From only
$200$ \emph{noisy} mixture experiments, both ensembles reach a cross-validated
$R^2\approx 0.56$ on final pH, and the strains they rank most important are exactly
the genuinely strong acidifiers: the random forest's top ten recover eight of the
ten strongest, and permutation importance correlates clearly with the true,
unseen, solo acidifying power. The covering array supplied the coverage; the trees
supplied the inference; together they turn an impossible combinatorial search into
a short list of strains worth pursuing.

\begin{figure}[t]
\centering
\includegraphics[width=\linewidth]{ferm_importance.png}
\caption[Tree-model recovery of influential strains]{Analysis of the $200$-run
screen. \emph{Left:} permutation importance, Eq.~\eqref{eq:permimp}, for the
top strains, from a random forest (blue) and XGBoost (orange); the two agree on the
leaders. \emph{Right:} each strain's recovered importance against its \emph{true}
acidifying power (negative solo final pH), which the screen never measured
directly---the strong positive trend shows the influential strains are recovered
from noisy blends alone.}
\label{fig:ferm-imp}
\end{figure}

\begin{prosbox}{tree-ensemble screening analysis}
\begin{itemize}
  \item \textbf{Interactions for free.} Trees capture threshold and synergy effects
  with no need to enumerate interaction terms---decisive when there are more
  possible pairs than runs.
  \item \textbf{Robust and honest.} Cross-validated $R^2$ reports true predictive
  power; permutation importance ranks drivers model-agnostically; agreement between
  forest and boosting guards against artefacts.
  \item \textbf{High-dimensional and forgiving.} They handle the wide, sparse,
  non-orthogonal presence matrix that defeats ordinary regression.
\end{itemize}
\end{prosbox}

\begin{consbox}{tree-ensemble screening analysis}
\begin{itemize}
  \item \textbf{Correlation, not mechanism.} Importance flags strains that
  \emph{predict} the outcome, not a causal pathway; confounded strains (often
  co-occurring by design) can share or steal credit.
  \item \textbf{Data-hungry for weak effects.} Subtle interactions need more than a
  few hundred runs to surface above the noise.
  \item \textbf{Less transparent than ANOVA.} There is no tidy table of signed
  coefficients with $p$-values---the deliverable a regulatory reviewer may expect.
\end{itemize}
\end{consbox}

%==============================================================================
\section{Synthesis}
\label{sec:ferm-synth}

\begin{keybox}{the living reactor}
A fermentation differs from a chromatography column in every respect that makes
modelling hard---it is alive, indirectly observed, intrinsically variable, and
chosen combinatorially---yet the same virtual-laboratory discipline applies. A
mechanistic--stochastic model (Sections~\ref{sec:ferm-model}--\ref{sec:ferm-noise})
supplies a faithful, randomness-aware oracle; a covering array
(Section~\ref{sec:ferm-screen}) interrogates a combinatorial space within a real
budget; tree-ensemble learning (Section~\ref{sec:ferm-analysis}) extracts the few
strains that matter from noisy data; and the fingerprint distance
(Section~\ref{sec:ferm-fingerprint}) turns ``reproduce the discontinued product''
into a scalar objective for the optimiser of Chapter~\ref{ch:design}. Model,
design, and analysis are, once again, three parts of one instrument---now pointed
at a beaker of milk.
\end{keybox}
